
%Muestra los márgenes del documento para evitar Warnings
%Para activar la siguiente línea quite el simbolo % 
%\usepackage[showframe]{geometry}

%Formato de fuentes bibliográficas
%Use el estilo bibliográfico que sea pertinente según el área de estudio APA, IEEE, etc

%Usando el paquete BibLaTeX
%Cita normal con \cite[página]{} y cita con paréntesis \parencite[página]{}

% Configuración de BibLaTeX
%\usepackage[backend=biber,style=authoryear,maxcitenames=2,maxbibnames=99,giveninits=true,uniquename=false]{biblatex}
%\addbibresource{biblio.bib}

% Cambiar el idioma de las referencias bibliográficas a español
%\DefineBibliographyStrings{spanish}{%
%  andothers = {et\addabbrvspace al\adddot},
%  andmore = {et\addabbrvspace al\adddot},
%}

% Personalizar el formato de las citas y la bibliografía
%\DeclareNameAlias{sortname}{family-given}
%\DeclareDelimFormat{multinamedelim}{\addcomma\space}
%\DeclareDelimFormat{finalnamedelim}{\addcomma\space\&\space}
%\DeclareFieldFormat{titlecase}{\MakeSentenceCase*{#1}}
%\DeclareFieldFormat[article,inbook,incollection,inproceedings,patent,thesis,unpublished]{title}{\titlecase{#1}}
%\DeclareFieldFormat{journaltitlecase}{\titlecase{#1}}
%\DeclareFieldFormat{pages}{#1}
%\DeclareFieldFormat{volume}{\mkbibbold{#1}}
%\renewbibmacro{in:}{}
%\AtEveryBibitem{\clearfield{month}}

%Usando el paquete Natbib
%Cita normal \cite[página]{} y cita con paréntesis \citep[página]{}
\usepackage[numbers,square]{natbib}
\bibliographystyle{dtvstyle}


%Idioma del documento
%Use main para el idioma principal del documento
\usepackage[main=spanish,english,german,french,portuguese]{babel}

% Carácteres especiales
\usepackage{fontenc}

% Evita ligadura li & fl
\usepackage{microtype}
\DisableLigatures{encoding = *, family = *}

\usepackage{amssymb}
% Otros paquetes de tablas y colores avanzados
\usepackage{amsmath,graphicx,rotating,float,multirow}
\usepackage{longtable}
\setlength{\LTcapwidth}{6in}
\usepackage[utf8]{inputenc}
\usepackage{epsfig,epic,eepic,threeparttable,amscd,here,lscape,tabularx,subfigure}
\usepackage{tabu,array}
\usepackage[rgb]{xcolor}

% Permite ver y configurar los parámetros de la página
\usepackage{layout}
%Hyperref permite ver las secciones del texto
\usepackage[hidelinks]{hyperref}

%Permite incluir código de cualquier lenguaje dentro del texto del documento
\usepackage{minted}
\usepackage{fancyvrb}
\newenvironment{myverbatim}{\Verbatim}{\endVerbatim}
 
\usepackage{enumitem}

%Genera los comandos de la página de autoría
\newcommand{\studentname}{Johan Alejandro López Arias}
\newcommand{\submissiondate}{}
\newcommand{\academictitle}{}
\newcommand{\resgroupone}{}
\newcommand{\resgrouptwo}{}
\newcommand{\researchtopic}{}
\newcommand{\thesisname}{}
\newcommand{\thesisnameeng}{}
\newcommand{\thesisnamelang}{} %Usar solo si se requiere
\newcommand{\director}{Pedro Fabián Cárdenas Herrera}
\newcommand{\directortitle}{Prof. Dr. Ing.}
\newcommand{\codirector}{Pedro Fabián Cárdenas Herrera} %Usar solo si se requiere
\newcommand{\codirectortitle}{} %Usar solo si se requiere
\newcommand{\issuedate}{}
\newcommand{\palabrasclave}{}
\newcommand{\keywords}{}
\newcommand{\schlusselworter}{}
\newcommand{\palavraschave}{}
\newcommand{\sede}{Bogotá}
\newcommand{\department}{Departamento de Ingeniería Mecánica y Mecatrónica}
\newcommand{\departmenttwo}{Departamento de Ingeniería Mecánica y Mecatrónica} %Usar solo si se requiere
\newcommand{\faculty}{}
\newcommand{\university}{Universidad Nacional de Colombia}

%Información de la tesis
%Diligenciar aquí los datos para su carga automática donde se requiera en el documento
%En el caso de tesis o trabajos finales, verificar que el título coincida con el aprobado por la Facultad
\renewcommand{\studentname}{Johan Alejandro López Arias}
\renewcommand{\thesisname}{Diseño de un Kit de Clasificación de Objetos con Visión por Computador para el Robot PhantomX Pincher Utilizando ROS2 y Raspberry Pi 5}
\renewcommand{\thesisnameeng}{Design of an Object Classification Kit with Computer Vision for the PhantomX Pincher Robot Using ROS2 and Raspberry Pi 5}
\renewcommand{\thesisnamelang}{Entwurf eines Objektklassifizierungs-Kits mit Computer Vision für den PhantomX Pincher Roboter unter Verwendung von ROS2 und Raspberry Pi 5} %Usar solo si se requiere
\renewcommand{\issuedate}{2025}
\renewcommand{\submissiondate}{2025}
\renewcommand{\director}{Prof. Dr. Ing. Pedro Fabián Cárdenas Herrera}
\renewcommand{\directortitle}{Profesor Titular}
\renewcommand{\codirector}{}
\renewcommand{\codirectortitle}{}
\renewcommand{\academictitle}{Ingeniero Mecatrónico - Modalidad Trabajo de Grado }
\renewcommand{\resgroupone}{Grupo A (Sigla Grupo Investigación 01) }
\renewcommand{\resgrouptwo}{Grupo B (Sigla Grupo Investigación 02) }
\renewcommand{\researchtopic}{Línea}
\renewcommand{\sede}{Sede Bogotá} 
\renewcommand{\department}{Departamento de Ingeniería Mecánica y Mecatrónica}
\renewcommand{\departmenttwo}{} %Usar solo si es necesario
\renewcommand{\faculty}{Facultad de Ingeniería}

%Palabras clave del documento - Tener presente los Theasurus https://www.thesaurus.com/
%Disponible en 3 idiomas aunque se puede extender a francés o otro idioma
\renewcommand{\palabrasclave}{ROS2, Visión por Computador, Clasificación de Objetos, Robótica Educativa, Sistemas Embebidos, Raspberry Pi}

\renewcommand{\keywords}{ROS2, Computer Vision, Object Classification, Educational Robotics, Embedded Systems, Raspberry Pi}

\renewcommand{\schlusselworter}{ROS2, Computersehen, Objektklassifizierung, Bildungsrobotik, Eingebettete Systeme, Raspberry Pi}
%\renewcommand{\palavraschave}{}

% Estilo de los encabezados y pies de página
\usepackage{fancyhdr}
\fancyhf{}%
\pagestyle{fancyplain}
\textheight22.5cm \topmargin0cm \textwidth16.5cm \headheight22pt
\oddsidemargin0.5cm \evensidemargin-0.5cm%
\fancypagestyle{plain}{
\fancyhead[RO,LE]{}
\fancyhead[RE,LO]{\small \textbf{\thesisname}}
\fancyfoot[CO,CE]{\thepage}
}
\pagestyle{fancy}
\fancyhf{}%
\renewcommand{\chaptermark}[1]{\markboth{\thechapter.\; #1}{}}
\renewcommand{\sectionmark}[1]{\markright{\thesection.\; #1}{}}
\fancyhead[LO,RE]{\leftmark}
\fancyhead[RO,LE]{\rightmark}
\fancyfoot[CO,CE]{\thepage}
\thispagestyle{fancy}%

\usepackage{titlesec}
% Permite personalizar los títulos de sección y de capítulos
% hang lo deja en el mismo renglón, display lo despliega
% Elimina el "Capitulo" y deja solo el número
\titleformat{\chapter}[hang]
  {\sffamily\Huge\bfseries}{\thechapter}{0.5cm}{\sffamily\Huge}
\titleformat{\section}[hang]{\sffamily\LARGE}{\thesection}{0.5cm}{}
\titleformat{\subsection}[hang]{\sffamily\Large}{\thesubsection}{0.5cm}{}
\titleformat{\subsubsection}[hang]{\sffamily\large}{\thesubsubsection}{0.5cm}{}
\titleformat{\paragraph}[runin]{\sffamily\normalsize}{}{}{\emph}

%Coloca anexo o apéndice en la Tabla de contenido
\usepackage[toc,page]{appendix}

% Configuración de las páginas en twoside-mode
% Permite ver y configurar los parámetros de la página
\setlength{\voffset}{-0.25in}
\setlength{\headwidth}{467pt}
\setlength{\headheight}{22pt}
\setlength{\oddsidemargin}{0pt}
\setlength{\evensidemargin}{0pt}
\setlength{\marginparwidth}{0pt}
\setlength{\marginparsep}{0pt}
\setlength{\parskip}{2em}
\setlength{\footskip}{20pt}
\setlength{\textheight}{650pt}
\setlength{\textwidth}{467pt}
\setlength{\headsep}{5pt}
\setlength{\parindent}{0pt}
\setlength{\baselineskip}{10pt plus 5pt minus 5pt}
\renewcommand{\theequation}{\thechapter-\arabic{equation}}
\renewcommand{\thefigure}{\textbf{\thechapter-\arabic{figure}}}
\renewcommand{\thetable}{\textbf{\thechapter-\arabic{table}}}

%Ajusta el espacio entre la etiqueta de figuras y tablas y su título en la lista de figuras y en la de tablas 
\usepackage{titletoc} 
\titlecontents{figure}[0em]{}{\thecontentslabel\hspace{1em}}{}{\titlerule*[1pc]{.}\contentspage}
\titlecontents{table}[0em]{}{\thecontentslabel\hspace{1em}}{}{\titlerule*[1pc]{.}\contentspage}

%Define la distancia de la primera linea de un parrafo a la margen
\parindent0cm 

%Espacio entre lineas
\renewcommand{\baselinestretch}{1}

%Permite personalizar el ajuste vertical mediante cajas
\usepackage{adjustbox}

%Para rotar texto, objetos, tablas y páginas.
\usepackage{rotating}

%Permite incluir mecanismos y reacciones químicas
\usepackage{tikz}
\usepackage{chemformula}
\usepackage{chemfig}

\usetikzlibrary{calc,arrows.meta,shapes.geometric,arrows,positioning}% per right to e left to
\tikzset{
myedge/.style={->, -{Latex[#1]}}
}

%Fuente de la presentación Ancizar Sans UNAL
%Para usar este compilado en Overleaf se debe usar el compilador XeLaTeX o LuaLaTeX!!
%Menu -> Compiler -> XeLaTeX o LuaLaTeX
%La siguiente línea debe comentarse si desea compilar con pdfLaTeX
%\RequireXeTeX

% Definición de la fuente Ancizar Sans
\newif\ifxetexorluatex

\ifxetexorluatex
  \usepackage{fontspec}
  \usefonttheme{serif}
  \setmainfont{AncizarSans}[Path=./AncizarSans/,Scale=1,Extension=.otf,UprightFont=*-Regular,BoldFont=*-Bold,ItalicFont=*-Italic,BoldItalicFont=*-BoldItalic]
\else
  % Si se compila con pdfLaTeX, cargar la fuente apropiada aquí
  \usepackage[T1]{fontenc}
\fi
% Metadatos del documento
\AtBeginDocument{%
	\hypersetup{
		pdfborder={0 0 0},
		pdfauthor={\studentname},
		pdfsubject={\thesisname}, 
		pdfcreator={\studentname},
		pdfproducer={\studentname},
	}
}

%Carga el simbolo de grado y el de Angstrom
\newcommand{\angstrom}{\textup{\AA}}
\newcommand{\grad}{$^{\circ}$}


%Nombres y formatos de títulos, tablas y figuras
%Use \sffamily para dejar con letra Sans Serif, sin etiqueta queda LaTeX clásico
\renewcommand{\listfigurename}{\sffamily Lista de figuras}
\renewcommand{\listtablename}{\sffamily Lista de tablas}
\renewcommand{\contentsname}{\sffamily Contenido}
\renewcommand{\chaptername}{\sffamily Capítulo}
\renewcommand{\tablename}{\scriptsize \centering \textbf{Tabla}}
\renewcommand{\figurename}{\scriptsize \centering \textbf{Figura}}
\renewcommand{\appendixname}{\sffamily Anexo}


